% Harvest 초안 마인드맵 — CVPR 양식(공식 author kit의 cvpr.sty), 내용은 논문이 아니라 마인드맵.
% 부분마다 "지금 생각 / 바뀐 점 / 방향 / [결정 필요]"만 적는다. 안 정해진 곳은 비워 둔다.
% Overleaf: 이 폴더의 main.tex, cvpr.sty, ieeenat_fullname.bst를 올리고 컴파일러는 XeLaTeX.
\documentclass[10pt,twocolumn,letterpaper]{article}
\usepackage[pagenumbers]{cvpr}   % 초안이라 review(익명·줄번호) 대신 쪽번호만
\usepackage{kotex}
\usepackage{enumitem}
\setlist{nosep,leftmargin=1.2em}
\definecolor{cvprblue}{rgb}{0.21,0.49,0.74}
\usepackage[breaklinks,colorlinks,allcolors=cvprblue]{hyperref}

% 빨간색: 사용자가 의도하고 말한 것 / 검정색: Claude의 조사·판단·제안
\newcommand{\intent}[1]{\textcolor{red}{#1}}
\newcommand{\idea}[1]{\textbf{지금 생각:} #1}
\newcommand{\chg}[1]{\textbf{바뀐 점:} #1}
\newcommand{\dirn}[1]{\textbf{방향:} #1}
\newcommand{\decide}[1]{\textbf{[결정 필요]} #1}

\def\paperID{*****}
\def\confName{CVPR}
\def\confYear{2027}

\title{Harvest --- 초안 마인드맵 (v3)}
\author{마지막 갱신: 2026-09-23 20:15 UTC \quad 초안 조사 단계, v3 원문 검증 반영}

\begin{document}
\maketitle

\begin{abstract}
(아직 비움. 컨트리뷰션이 정해지면 쓴다.)
\end{abstract}

\section{Introduction --- 큰 줄기}
\begin{itemize}
  \item \intent{기존 스킬 기반 잡기를 기본으로 쓰고, 실패할 때마다 GPT-6 Astra가 개입해서 이어감}
  \item \intent{상위 Astra(느림, 계획) / 하위 Jev(빠른 객관식 판단). 예전 VLA의 큰 LLM + 작은 VLM처럼}
  \item \intent{VLA는 다른 환경으로 일반화가 안 됨. LLM은 됨. 항상 부각}
  \item \intent{안전은 크게 다루지 않음}
  \item \idea{계층 구조, Astra로 로봇 제어, 스텝마다 LLM 이산 행동은 이미 있음. 우리 자리는 Jev 확률로 겹쳐 부르고 스스로 확인하는 부분(M4)}
  \item \decide{주장 범위. $\pi_{0.5}$ 등이 새 환경을 주장함 $\rightarrow$ ``대규모 로봇 데이터·학습 없이 덜 무너짐''으로 좁힐지}
\end{itemize}

\paragraph{컨트리뷰션 후보 (v3)}
\begin{enumerate}
  \item M4: Jev 확률 기반 계단식 겹침 호출 + 측정 상태로 겹침 반영 확인·갱신 (선점 없음, 가장 강함)
  \item M8: 비정지 로봇에서 호출 시점 $\times$ effort 평가 (방법 아닌 평가 컨트리뷰션으로 좁힘)
  \item M6: 스킬 내부 결정 지점을 typed 질문으로
  \item M10: Astra가 컴파일한 규칙만 Jev에 (로봇 미적용)
\end{enumerate}
\chg{v2의 ``호출 방식 비교가 새롭다''는 틀림. 주기 대 이벤트 비교는 이미 여러 편}

\section{Related Work --- 선점 지도}
\begin{itemize}
  \item Astra로 로봇: AGP, GPT-as-Policy, TGL, RoboDojo (모두 미심사), Harness VLA가 Astra 계획기 사용
  \item 스텝마다 LLM이 이산 행동: Show-Harness $\rightarrow$ M3 직접 선행, 반드시 인용
  \item 주기 대 이벤트 호출: CheckVLA, BRACE(ICML 2026) 등
  \item Jev를 로봇에 쓴 논문은 아직 없음. 저장소 데모만(Jev-as-Policy 감시 중)
\end{itemize}

\section{Method --- 모듈별 생각}

\subsection*{두 모델 제약}
\begin{itemize}
  \item Jev: 텍스트만. 수치·시간 비교 약함. 입력 길면 정확도 하락. \textbf{Score 기댓값으로 크기 보간 금지}, \textbf{질문끼리 확률이 서로 안 맞음}(공식 문서)
  \item Astra: 이미지 OK, 비디오 없음. \textbf{logprob 못 씀}(effort none 없음). effort 지연은 날마다 바뀜
\end{itemize}

\subsection*{M1 이미지 $\rightarrow$ Jev}
\begin{itemize}
  \item \intent{Jev는 이미지를 못 봄. AI 전 분야에서 가장 효율적인 변환을 가져옴}
  \item \idea{형식보다 내용. 작업 관련 물체만 + 코드가 계산한 관계·술어}
  \item \chg{``텍스트가 이미지보다 낫다''의 근거(BALROG)는 모델마다 다름. 근거에서 내림}
  \item \decide{변환 방법, 인식 앞단}
\end{itemize}

\subsection*{M2 Astra $\rightarrow$ Jev}
\begin{itemize}
  \item \intent{M1과 같은 문제}
  \item \idea{물체 번호로 말하게 함(텍스트판 Set-of-Mark). 제약은 Astra가 코드로 쓰고 Jev에는 참/거짓만}
\end{itemize}

\subsection*{M3 행동}
\begin{itemize}
  \item \intent{촘촘하게 가도 됨}
  \item \chg{``확률 평균으로 크기''는 Jev 공식 문서가 금지 $\rightarrow$ 이름 붙은 구간을 고름}
  \item \idea{보기 255개 한 번이면 충분할 수 있음(해상도 250 대 1000 차이 작음)}
  \item \dirn{첫 실험 = Jev 보정 측정}
\end{itemize}

\subsection*{M4 RTC식 겹침 호출 ★}
\begin{itemize}
  \item \intent{1초에 여러 번 계단식으로 겹쳐 호출}
  \item \intent{겹친 구간이 제대로 반영됐는지 스스로 확인하고 업데이트. 가장 중요}
  \item \idea{확인을 둘로 나눔: (a) 호출끼리 같은 말을 하나(앞부분 합의) (b) 실제로 그렇게 움직였나(예상 상태 대 측정 상태, 코드)}
  \item \idea{(a)만으론 안 됨. 같은 모델은 같은 오답을 반복}
  \item \chg{VLASH·A2C2는 학습형이라 발상만. A2C2는 ICLR 탈락}
\end{itemize}

\subsection*{M5 스무딩}
\begin{itemize}
  \item \intent{아주 작은 움직임은 스무딩으로. 두 단계}
  \item \idea{v2 그대로(겹침 블렌딩 + jerk 최소화)}
\end{itemize}

\subsection*{M6 스킬 $\leftrightarrow$ Jev}
\begin{itemize}
  \item \intent{잘게 엮음. 1년 이내 스타 상위 5개 참고}
  \item \idea{스타 상위: PhyAgentOS, Zetta, Harness VLA, CaP-X, ROSClaw(또는 Show-Harness)}
  \item \idea{스타 1위는 거칠게 엮음. 잘게 엮는 쪽은 Zetta(critic 코드 + 수락/거부)와 Show-Harness}
  \item \decide{5위 ROSClaw 대 Show-Harness}
\end{itemize}

\subsection*{M7 진행 판별}
\begin{itemize}
  \item \intent{VLA의 ``계획대로'' 정의를 참고}
  \item \idea{v2 그대로(센서·술어, 남은 시간)}
\end{itemize}

\subsection*{M8 Astra 호출}
\begin{itemize}
  \item \intent{실패 판정 때 부름. 연속 프레임(예: 한 이미지에 10개)}
  \item \intent{비교: 일정 주기 / 정한 순간 / ``이때쯤 나왔어야 하는데''}
  \item \intent{처음 계획만 멈춤}
  \item \idea{``이때쯤''은 코드가 예상 시간으로 계산. LLM에게 맡기면 안 됨}
  \item \idea{프레임은 격자보다 ``어떤 프레임을 고르냐''가 중요. 움직임 큰 순간 위주 6--10장}
  \item \chg{effort low 기본안 삭제(사용자: 추론 강도 낮추지 말 것). effort는 실험 축만}
\end{itemize}

\subsection*{M9 실패 복구}
\begin{itemize}
  \item \intent{VLA에서 가장 좋은 실패 복구. 1년 반 이내}
  \item \idea{아래 층부터: 스킬 재시도 $\rightarrow$ Jev $\rightarrow$ 리셋 스킬 $\rightarrow$ Astra}
  \item \chg{FLARE는 자세 오류/환경 파괴 구분만 가져옴(재시도는 학습형이었음)}
\end{itemize}

\subsection*{M10 경험 축적}
\begin{itemize}
  \item \intent{API 모델 기반 경험 축적을 제대로}
  \item \idea{Astra: 단순 경험 검색(ExpRAG) + 교훈. Jev: 교훈 목록 말고 딱 맞는 규칙 0--2개만}
  \item \chg{``메모리 없는 게 낫다''는 반대 논문이 실재(EMNLP 2026) $\rightarrow$ 같은 예산 비교 필수}
\end{itemize}

\section{Experiments --- 평가 방향}
\begin{itemize}
  \item \intent{비교: Astra만 / Jev만 / VLA / 룰베이스 / LLM+스킬 기존 논문 / Astra 쓰는 기존 논문}
  \item \idea{``Astra 쓰는 기존 논문''이 이제 있음 $\rightarrow$ 재현 가능한 것 선택}
  \item \idea{AGNOSTOS의 LLM 방법은 정답 좌표를 받음 $\rightarrow$ 정답 상태 / 인식 상태 두 조건}
  \item \decide{실험 환경}
\end{itemize}

\end{document}
